% A new mathematical supplement. No claim of checker-size improvement.
% Requires amsmath, amssymb, amsthm; use own theorem environments if needed.
\section{Finite scalar clocks without induction}
\label{sec:finite-clocks}

This section resolves the even-cycle gap in the linked README's first
further question. It does not remove the progress rule or establish a smaller
complete checker. The result is a constraint on any proposed countermodel
for the independence of progress.

Work first in intuitionistic first-order logic with actual equality, the
two scalar iterator computations, and numerical coverage
\[
 x=0\ \lor\ \exists y\,x=Sy.
\]
There is no induction rule. Write
\[
 C(x,a,b)=I(x,a;u.b),\qquad Z(x)=C(x,1,0),\qquad
 \bar j=S^j0.
\]
Thus $C(0,a,b)=a$, $C(Sx,a,b)=b$, and coverage proves
$Z(x)=0\lor Z(x)=1$. All inductions below are external inductions on a
fixed finite integer specifying a term, not induction in the object theory.
For a fixed unary term $h$, $h^{[r]}$ means $r$-fold syntactic composition;
it does not introduce a new iterator or an object-language quantifier.

\begin{lemma}[Finite clocks]
For every fixed integer $m\geq2$, there is a finite scalar term $h_m(x)$
such that the theory without induction proves
\[
 h_m(x)\in\{\bar0,\ldots,\overline{m-1}\},\qquad
 h_m(\bar j)=\overline{j+1}\quad(0\leq j<m-1),\qquad
 h_m(\overline{m-1})=0.
\]
Here membership abbreviates a finite disjunction of actual equalities.
\end{lemma}
\begin{proof}
Start with $h_2=Z$. Given $h_m$, define
\[
 T_m(x)=I(x,0;u.h_m(u)).
\]
Coverage on $x$ and the two computations show that $T_m(x)$ belongs to
$\{\bar0,\ldots,\overline{m-1}\}$: the zero case is $0$, and in the
successor case its value is an $h_m$-value. No assertion about the preceding
iterator value is needed.

For $1\leq r<m$, put
\[
 d_{m,r}(z)=Z\bigl(h_m^{[m-r]}(z)\bigr),\qquad
 V_{m,0}(z)=0,
\]
\[
 V_{m,r}(z)=C\bigl(d_{m,r}(z),V_{m,r-1}(z),\overline{r+1}\bigr),
 \qquad V_m=V_{m,m-1}.
\]
Finite ground computation proves
$d_{m,r}(\bar s)=1$ for $s=r$, and $d_{m,r}(\bar s)=0$ for
$0\leq s<m$, $s\ne r$. Consequently
\[
 V_m(0)=0,\qquad V_m(\bar r)=\overline{r+1}\quad(1\leq r<m).
\]
Every value of $V_m$, at an arbitrary argument, is one of
$0,\bar2,\ldots,\bar m$, by finitely many applications of coverage to
the selector counters. Define
\[
 \boxed{\quad h_{m+1}(x)=C\bigl(x,1,V_m(T_m(x))\bigr).\quad}
\]
Its range is $\{0,\bar1,\ldots,\bar m\}$. At $x=0$ its value is $1$.
At $x=\bar j$ with $1\leq j<m$, finite computation gives
$T_m(\bar j)=\bar j$, so its value is $\overline{j+1}$. At $x=\bar m$,
$T_m(\bar m)=0$, so its value is $0$. This completes the external
construction.
\end{proof}

The first nontrivial extension has the compact description
\[
 h_3(x)=C\bigl(x,1,C(Z(Z(T_2(x))),0,2)\bigr),
 \qquad T_2(x)=I(x,0;u.Z(u)).
\]
Direct computation gives $h_3(0)=1$, $h_3(1)=2$, and $h_3(2)=0$.

\begin{lemma}[Separation of finitely many numerals]
For fixed $m\geq2$ and distinct $i,j<m$, the assumption
$\bar i=\bar j$ proves $0=1$, without induction or successor injectivity.
\end{lemma}
\begin{proof}
Apply $h_m^{[r]}$ with $r\equiv-i\pmod m$ and $0\leq r<m$.
The resulting equation is $0=\bar s$ for some $1\leq s<m$.
Applying $Z$ gives $1=0$, whose symmetric form is the desired
contradiction. Only finitely many computation equations are used.
\end{proof}

\begin{theorem}[No finite successor cycle without induction]
For every fixed integer $k\geq1$ and every scalar term $c$, the theory
without induction derives
\[
 S^k c=c\ \longrightarrow\ 0=1.
\]
With $0\ne1$ it therefore proves $S^k c\ne c$.
\end{theorem}
\begin{proof}
Choose $m=k+1$ and set $v=T_m(c)$. The hypothesis and $k$ applications of
successor computation give
\[
 v=T_m(c)=T_m(S^k c)=h_m^{[k]}(v).
\]
The finite-range property of $T_m$ gives the cases $v=\bar j$ for
$0\leq j<m$. In each case the displayed equation becomes
\[
 \bar j=\overline{j+k\pmod m}.
\]
Since $k=m-1$, the two residues differ. The preceding separation lemma
gives $0=1$ in every case.
\end{proof}

\subsection{Removing the finite disjunctions in the native calculus}

There is also a direct written translation into the native kernel without
progress, so the preceding use of ordinary disjunction is dispensable.
The existing theorem \texttt{NoInd.cases} derives the following rule for
equation templates: if $E(0)$ and $E(Sz)$ are derivable uniformly in the
retained context, then $E(t)$ is derivable for every term $t$. The
computation and crossed equality rules provide ordinary equational
congruence and transitivity without progress.

Call a term a finite selector tree over $L$ if it is built from leaves in
the fixed finite list $L$ by $C(q,a,b)$, with arbitrary counter terms $q$.
For any equation template $E$, if $E(\ell)$ is derivable uniformly for
each leaf $\ell\in L$, then $E(t)$ is derivable for every such tree $t$.
Prove this externally by induction on its finite syntax: at a selector,
apply \texttt{NoInd.cases} to its counter and compute the two branches.
Each $h_m(x)$ is literally a finite selector tree over
$L_m=\{\bar0,\ldots,\overline{m-1}\}$. A further application of
\texttt{NoInd.cases} to the counter of $T_m(c)$ shows:
\[
 \bigl(E(\bar j)\text{ derivable for every }j<m\bigr)
 \quad\Longrightarrow\quad E(T_m(c))\text{ derivable}.
 \tag{FR}
\]
The zero branch is $E(0)$; the successor branch is
$E(h_m(T_m(z)))$, handled by the selector-tree argument.

Put $b_r(a)=Z(h_m^{[(-r\bmod m)]}(a))$ for $0\leq r<m$.
These are Boolean terms, with $b_r(\bar j)=1$ exactly when $r=j$.
For Boolean arguments define
\[
 \operatorname{xor}(p,q)=C(p,q,Z(q)),\qquad
 \operatorname{or}(p,q)=C(p,q,1).
\]
Let $D_m(a,b)$ be the left-associated finite disjunction (starting at
$0$) of the terms $\operatorname{xor}(b_r(a),b_r(b))$.
Finite zero/successor cases prove $D_m(a,a)=0$ for every $a$; no
disjunction formula or equality characteristic for arbitrary numbers is
used. For $k=m-1$, ground computation proves
\[
 D_m(\bar j,h_m^{[k]}(\bar j))=1\qquad(j<m).
\]
Apply (FR) to this equation template to obtain
$D_m(v,h_m^{[k]}(v))=1$ at $v=T_m(c)$. The cycle assumption gives
$v=h_m^{[k]}(v)$, so equational transport and $D_m(v,v)=0$ give
$0=1$. This is the required native derivation, with no progress rule.

\paragraph{Boundary of the result.}
The construction and the native translation are written proofs; no new
Lean verification is claimed. They do not prove the progress rule
redundant, and do not by themselves settle whether successor uniqueness
is equivalent to full iPRA. The novelty is the general finite-clock
construction, which completes the existing odd-cycle exclusion.

The companion \texttt{check\_finite\_clocks.py} builds genuine scalar
terms using de Bruijn variables and evaluates them on bounded standard
inputs. It stores repeated subexpressions as a directed acyclic graph;
this is only a representation of a finite expanded term. It represents
$h_m^{[r]}(a)$ as $I(\bar r,a;u.h_m(u))$, which expands to the same
composition by finitely many computations. Its default run checks
periods $2$ through $12$ and numerical inputs $0$ through $32$, including
finite range, clock values, cycle transitions, and the Boolean
discriminator: 1,243 checks pass. These evaluations corroborate finite
examples only; the external construction and proofs above establish the
universal theorem.
